\scanpage{197}
\begin{center}- \mspage{59} -\end{center}
\margin{IX}

désignant respectivement les parties abéliennes de $A_{\unsure{o|0}}^{\unsure{0|o}}$ et $A'_{\unsure{o|0}}{}^{\unsure{0|o}}$. Comme, en vertu de VIII 4.8, la biextension $W_{\unsure{o|0}}$ provient d'une biextension $W_{\unsure{o|0}}^{\mathrm{ab}}$ de $(B_{\unsure{o|0}},B'_{\unsure{o|0}})$ par $\mathbb G_{mk}$, déterminée à isomorphisme unique près, l'accouplement (5.3.5) peut aussi s'interpréter comme déduit de la correspondance divisorielle $W_{\unsure{o|0}}^{\mathrm{ab}}$, et l'homomorphisme correspondant
\[
(5.3.6)\qquad T_\ell(B_{\unsure{o|0}})\longrightarrow D(T_\ell(B'_{\unsure{o|0}}))\simeq T_\ell(D(B'_{\unsure{o|0}}))
\]
(où dans le dernier terme, la lettre $D$ désigne le schéma abélien dual) n'est autre que celui déduit par fonctorialité de $T_\ell$ de l'homomorphisme
\[
(5.3.7)\qquad B_{\unsure{o|0}}\longrightarrow D(B'_{\unsure{o|0}})
\]
défini par $W_{\unsure{o|0}}^{\mathrm{ab}}$. Donc le fait que (5.3.5) soit une dualité parfaite signifie aussi, par le dictionnaire bien connu (cf. [ 5]) que le noyau de (5.3.7), qui est un groupe fini sur $k$ en vertu de 1.5, a une composante $\ell$-primaire nulle. Dire qu'il en est ainsi pour tout $\ell$ signifie donc que l'isogénie (5.3.7) est un isomorphisme, i.e. que la correspondance divisorielle $W_{\unsure{o|0}}^{\mathrm{ab}}$ est une dualité parfaite entre $B_{\unsure{o|0}}$ et $B'_{\unsure{o|0}}$. Les considérations qui précédaient montrent seulement que le noyau de (5.3.7) est un $p$-groupe. Enonçons donc le résultat plus complet :

\uline{Théorème} 5.4. \uline{Soit} $S$ \uline{un trait}, \uline{et} $A_K$ \uline{un schéma abélien sur} $S$ \uline{ayant une réduction semi-stable sur} $S$. \uline{Désignons par} $B_{\unsure{o|0}},B'_{\unsure{o|0}}$ \uline{les parties abéliennes des modèles de Néron de} $A_K$ \uline{et du schéma abélien dual} $A'_K$ \uline{respectivement}. \uline{Alors la correspondance divisorielle canonique} $W_{\unsure{o|0}}^{\mathrm{ab}}$ \uline{sur} $(B_{\unsure{o|0}},B'_{\unsure{o|0}})$, \uline{déduite par} VIII 4.8 \uline{de la biextension canonique} $W_{\unsure{o|0}}$ \uline{de} $(A_{\unsure{o|0}}^{\unsure{0|o}},A'_{\unsure{o|0}}{}^{\unsure{0|o}})$ (\uline{fibre spéciale de la biextension} $W$ \uline{de} 1.4), \uline{est une dualité parfaite entre} $B_{\unsure{o|0}}$ \uline{et} $B'_{\unsure{o|0}}$,
\margin{377}
\margin{377}
\margin{196}
